\section{加入者}

\subsection{加入者の範囲と一定の資格}

\begin{InyouJobunbox}
  \begin{itemize}
    \item 法第25条
    \item 解釈第1-1
  \end{itemize}
\end{InyouJobunbox}

\begin{fillblank}[数理人試験 \ H28 \ 問題1 \ 設問2]
  \noindent \textbf{◯法令解釈}
  \begin{lawstructure}
    \begin{lawarticle}{}{}
      \lawparagraph{1} 加入者とすることについての「一定の資格」の内容
      
      （略）

      \lawitem{(1)}
      「一定の資格」として定めることができる資格とは、次の\textcircled{1}から\textcircled{4}に掲げる資
      格であり、これら以外のものを「一定の資格」として定めることは、基本的には
      特定の者に不当に差別的な取扱いとなるものであること。

      \hangedlawitem{\textcircled{1}}{3}
      「一定の職種」\\   \hspace*{1.0em}%
        「一定の職種」に属する従業員のみを加入者とすること。この場合におい
        て、「職種」とは、研究職、営業職、事務職などの労働協約又は就業規則その
        他これらに準ずるもの（以下「労働協約等」という。）において規定される職
        種をいい、これらの職種に属する従業員に係る給与及び退職金等の労働条件が
        他の職種に属する従業員の労働条件とは別に規定されているものであること。

      \hangedlawitem{\textcircled{2}}{3}
      「一定の（　　A　　）」、「一定の年齢」\\      \hspace*{1.0em}%
      従業員が労働協約等に定める見習期間中若しくは試用期間中であること又は
      労働協約等に定める退職金の算定対象期間に含まれていない期間中であること
      など加入者としないことに合理的な理由がある場合にあって、「一定の（　　A　　）
      」以上又は「一定の年齢」以上若しくは以下の従業員のみを加入者とするこ
      と。この場合にあっては、「一定の（　　A　　）」以上の従業員のみを加入者とす
      る場合にあっては（　　B　　）以上の（　　A　　）を有する従業員について、「一定の年
      齢」以上の従業員のみを加入者とする場合にあっては（　　C　　）以上の従業員につ
      いて、「一定の年齢」未満の従業員のみを加入者とする場合にあっては（　　D　　）
      未満の従業員については、少なくともこれを加入者とするものであること。     

      \hangedlawitem{\textcircled{3}}{3}
      「希望する者」\\      \hspace*{1.0em}%
      従業員のうち、「加入者となることを希望した者」のみを加入者とすること
      （この場合にあっては、加入者がその資格を（　　E　　）することを任意に選択できる
      ものではなく、かつ、将来にわたって安定的な（　　F　　）が確保されるように制
      度設計上配慮されていること。また、令第４条第１号の規定により、加入者は
      その資格を任意に（　　E　　）することはできないこととされていること。）。

      \hangedlawitem{\textcircled{4}}{3}
      「休職等期間中ではない者」\\      \hspace*{1.0em}%
      従業員のうち、「休職等（労働協約等に規定される育児休業、介護休業を含
      む。以下同じ。）期間中の者」について、「休職等期間の全部又は一部」が労
      働協約等に定める退職金の算定対象期間に含まれていないなど加入者としない
      ことに合理的な理由がある場合にあっては、当該「休職等期間の全部又は一
      部」に該当する従業員以外の従業員のみを加入者とすること。

    \end{lawarticle}
  \end{lawstructure}
\end{fillblank}

\begin{answerbox}
  \begin{enumerate}[A:]
    \item 勤続期間
    \item 5年
    \item 30歳
    \item 50歳
    \item 喪失
    \item 加入者数
  \end{enumerate}
\end{answerbox}

\begin{shortanswer}[年金1 \ 2023 \ 問題2(1)ウ]
  加入者の資格を区分(グループ区分)することができる場合として、『確定給付企業年金法並び
にこれに基づく政令及び省令について（法令解釈）』に定める場合を簡潔に入力しなさい。
（250 字以内）
\end{shortanswer}

\begin{answerbox}
  労働協約等における給与や退職金等の労働条件が異なるなど、合理的な理由がある場合
\end{answerbox}

\begin{explanationbox}
  法令解釈1-1(3)
\end{explanationbox}

\begin{truefalse}[年金1 2020 問題1(3)\textcircled{3}]
\textbf{下線部の内容の正誤を判定し、誤っている場合は正しい内容に改めなさい。}
\\

  実施事業所に使用される厚生年金保険の被保険者が加入者となることについて一定の資格を
定める場合、「一定の資格」として定めることができる資格とは、 \underline{「一定の職種」、「一定の勤続
期間」、「希望する者」}であり、これら以外のものを「一定の資格」として定めることは、基本
的には特定の者に不当に差別的な取扱いとなる。
\end{truefalse}

\begin{answerbox}
  \textbf{誤}

  「一定の職種」、「一定の勤続年数」、「一定の年齢」、「希望する者」、「休職等期間等ではない者」
\end{answerbox}

\begin{truefalse}[年金1 2018 問題1(1)\textcircled{2}]
  \textbf{下線部の内容の正誤を判定し、誤っている場合は正しい内容に改めなさい。}
  \\

  一定の年齢以上の従業員のみを加入者とする場合にあっては \underline{30 歳}以上の従業員について
は、少なくとも加入者としなければならない。
\end{truefalse}

\begin{answerbox}
  \textbf{正}
\end{answerbox}

\begin{truefalse}[年金1 2021 問題1(3)\textcircled{2}]
  \textbf{下線部の内容の正誤を判定し、誤っている場合は正しい内容に改めなさい。}
  \\

  規約において加入者の資格を区分しており、それらの区分（ 「区分Ａ」および「区分Ｂ」とす
る）の間で支給要件並びに給付の額の算定方法が異なっている場合において、加入者の資格を
取得したときに区分Ａだった加入者が、加入者の資格を喪失したときに区分Ｂとなっていたと
きには、\underline{区分Ｂにおける支給要件を適用するとともに、資格喪失} \\
\underline{時に区分Ａであったとした場合の給付の額を保証すること}とされている。
\end{truefalse}

\begin{answerbox}
  \textbf{誤}

  区分Aにおける支給要件を適用し（加入者の利益になる場合等を除く）、区分Bへ変更する前日に加入者の資格を喪失したとした場合の給付の額を保証すること。
\end{answerbox}

\begin{explanationbox}
  参照：規約例規第３条留意事項
\end{explanationbox}

\begin{truefalse}[年金1 2022 問題1(3)ウ]

    \textbf{下線部の内容の正誤を判定し、誤っている場合は正しい内容に改めなさい。}
  \\

  確定給付企業年金法第4 条第4 号における加入者となることについての「一定の資格」とし
て「希望する者」を定める場合は、加入者がその資格を喪失することを任意に選択できるも
のではなく、かつ、\underline{財政の健全性}が確保されるように制度設計上配慮されている必要があ
る。
\end{truefalse}

\begin{answerbox}
    \textbf{誤}

    将来にわたって安定的な加入者数
\end{answerbox}

\begin{shortanswer}[年金1 2025 問題2(1)エ]
  一定の勤続期間以上または一定の年齢以上若しくは一定の年齢以下の従業員のみを加入者とす
る場合において、少なくとも加入者としなければならない従業員の範囲として通知「確定給付
企業年金法並びにこれに基づく政令及び省令について（法令解釈）」 に定められている内容を簡
潔に入力しなさい。 （250 字以内）
\end{shortanswer}

\begin{answerbox}
  以下の従業員は少なくとも加入者にしなければならない。
  \begin{itemize}
    \item 一定の勤続期間以上のみ：5年以上の勤続年数を有する従業員
    \item 一定の年齢以上のみ：30歳以上の従業員
    \item 一定の年齢未満のみ：50歳未満の従業員
  \end{itemize}
\end{answerbox}

\newpage

\subsection{資格取得の時期}

\begin{InyouJobunbox}
  \begin{itemize}
    \item 法第26条、第103条
  \end{itemize}
\end{InyouJobunbox}

\begin{shortanswer}[自作 2026/10/11]
  確定給付企業年金法第26条に規定されている加入者の資格の取得時期について簡記せよ。
\end{shortanswer}

\begin{answerbox}
  以下のいずれかに該当するに至ったとき、加入者の資格を取得する。
  \begin{itemize}
    \item 実施事業所に使用されるに至ったとき
    \item その使用される事業所若しくは事務所又は船舶が、実施事業所となったとき
    \item 実施事業所に使用される者が、厚生年金保険の被保険者となったとき
    \item 実施事業所に使用される者が、規約に定められている資格を取得したとき
  \end{itemize}
\end{answerbox}

\begin{truefalse}[自作 2026/10/11]
    \textbf{下線部の内容の正誤を判定し、誤っている場合は正しい内容に改めなさい。}
  \\

  規約上、「30歳に達した日」に加入者資格を取得すると規定されている場合、
  この「30歳に達した日」とは\underline{「30歳の誕生日」}を指している。
\end{truefalse}

\begin{answerbox}
  \textbf{誤}

  30歳の誕生日の前日
\end{answerbox}

\begin{explanationbox}
  規約例規第４条代替例４

  参考：年齢計算ニ関スル法律
\end{explanationbox}

\newpage

\subsection{資格喪失の時期}

\begin{InyouJobunbox}
  \begin{itemize}
    \item 法第27条
  \end{itemize}
\end{InyouJobunbox}

\begin{shortanswer}[自作 2026/10/11]
  確定給付企業年金法第27条に規定されている加入者の資格喪失時期について簡記せよ。
\end{shortanswer}

\begin{answerbox}
  以下のいずれかに該当するに至ったとき、加入者の資格を喪失する。
  \begin{itemize}
    \item 死亡したとき
    \item 実施事業所に使用されなくなったとき
    \item その使用される事業所又は船舶が、実施事業所でなくなったとき
    \item 厚生年金保険の被保険者でなくなったとき
    \item 規約に定められている資格を喪失したとき
  \end{itemize}  
\end{answerbox}

\newpage

\subsection{加入者期間}

\begin{InyouJobunbox}
  \begin{itemize}
    \item 法第28条
    \item 令第21条、第22条
  \end{itemize}
\end{InyouJobunbox}

\begin{truefalse}[年金1 2022 問題1(3)イ]
ある加入者について、給与比例年金部分と定額年金部分の両方の給付がある制度において、
支給要件の判定に用いる加入者期間がそれぞれの年金部分ごとに異なることは、\underline{労働協約等
に定められている場合}など特別の事情によるやむを得ない場合を除き認められない。
\end{truefalse}

\begin{answerbox}
  \textbf{誤}

  他制度から権利義務を承継した場合
\end{answerbox}

\begin{explanationbox}
  規約例規第6条、Q\&A項番10

  年金部分ごとに加入者期間の端数処理を変更することも認められない。
\end{explanationbox}

\begin{shortanswer}[年金1 2020 問題2(1)\textcircled{1}]
  確定給付企業年金法施行令第 22 条第1 項に定める、加入者期間に算入することができる
「加入者となる前の期間」について簡記しなさい。
\end{shortanswer}

\begin{answerbox}
  \begin{enumerate}
    \item 当該確定給付企業年金の加入者に係る確定給付企業年金の実施前の期間のうち当該確定給付企業年金が実施されていたとしたならばその者が加入者となっていたと認められる期間その他これに準ずる期間
    \item 当該確定給付企業年金の加入者の資格を取得する前にその実施事業所に使用されていた期間の全部または一部
    \item 他の厚生年金適用事業所に使用されていた期間の全部または一部（規約において当該他の厚生年金適用事業所の名称および所在地ならびに加入者期間に算入する期間が定められている場合に限る。）
  \end{enumerate}
\end{answerbox}

\begin{fillblank}[自作 2026/10/11]
  \noindent \textbf{◯DB令}
  \begin{lawstructure}
    \begin{lawarticle}{第二十一条}{再加入者の加入者期間の合算に関する基準}
      法第二十八条第二項の政令で定める基準は、加入者の資格を喪失した後、再びもとの確定給付企業年金の加入者の資格を取得した者（以下「再加入者」という。）のうち、次に掲げるものについては、当該確定給付企業年金における前後の加入者である期間（以下「加入者期間」という。）を合算しないものであることとする。

      \lawitem{一}
      再加入者となる前に当該確定給付企業年金の（　　A　　）（法第二十九条第一項第二号に規定する（　　A　　）をいう。以下同じ。）の受給権者（給付を受ける権利（以下「受給権」という。）を有する者をいう。以下同じ。）となった者であって当該（　　A　　）の全部を支給されたもの（当該再加入者となったときに当該確定給付企業年金の（　　B　　）（同条第二項第一号に規定する（　　B　　）をいう。以下同じ。）の受給権者である者を除く。）

      \lawitem{二}
      再加入者となる前に当該確定給付企業年金の（　　C　　）（法第二十九条第一項第一号に規定する（　　C　　）をいう。以下同じ。）の受給権者となった者であって当該（　　C　　）の全部を支給されたもの（当該再加入者となったときに当該確定給付企業年金の（　　B　　）の受給権者である者を除く。）

      \lawitem{三}
      再加入者となる前に当該確定給付企業年金の（　　B　　）の受給権者となった者であって当該（　　B　　）の全部を支給されたもの（当該再加入者となったときに当該確定給付企業年金の（　　C　　）又は（　　A　　）の受給権者である者を除く。）

      \lawitem{四}
      加入者の資格を喪失した後に法第八十一条の二第二項、第八十二条の三第二項又は第九十一条の十九第二項の規定により（　　A　　）相当額が移換された者
    \end{lawarticle}
  \end{lawstructure}
\end{fillblank}

\begin{answerbox}
  \begin{enumerate}[A:]
    \item 脱退一時金
    \item 障害給付金
    \item 老齢給付金
  \end{enumerate}
\end{answerbox}

\newpage

\subsection{資格取得・喪失に関する手続き}

\begin{InyouJobunbox}
  \begin{itemize}
    \item 則第22条〜第23条の3
  \end{itemize}
\end{InyouJobunbox}

\begin{truefalse}[自作 2026/10/11]
    \textbf{下線部の内容の正誤を判定し、誤っている場合は正しい内容に改めなさい。}
  \\
  
  基金型DBの事業主は、その使用する者が基金の加入者の資格を取得または喪失したとき、
  その資格の取得または喪失した日から起算して\underline{30日を経過する日まで}に
  DB則第22条及び第23条の事項を基金に届け出なければならない。
\end{truefalse}

\begin{answerbox}
  \textbf{誤}

  30日を経過する日又は当該資格を取得又は喪失した日の属する月の翌月14日のいずれか早い日まで
\end{answerbox}

\begin{explanationbox}
  毎月末日現在における他制度加入者に関する他制度掛金相当額を翌月末日までに企業年金連合会を経由して国民年金基金連合会に通知する必要があるため、2024年12月1日改正で従前よりも早くなった。
\end{explanationbox}

\newpage